\section{The Roman College and the Making of the \emph{Disputationes}}

\subsection{A chair for contested doctrine}

In 1576 Bellarmine returned to Rome to occupy the Roman College chair of
Controversies.  He taught there into the later 1580s; institutional summaries
differ slightly about the endpoint.  The chair embodied a post-Tridentine
priority.  Clergy and teachers needed to know not only Catholic conclusions
but the scriptural, historical, and ecclesiological arguments by which
conflicting claims could be assessed.

Bellarmine's lectures brought a wide field into order.  Scripture, councils,
Fathers, scholastic analysis, and the texts of adversaries were coordinated
around disputed questions.  This created an unusually portable Catholic
argument: a student could see the structure of a controversy and the place of
particular authorities within it.  The method's reach should not be confused
with neutrality.  These were Catholic lectures designed to refute positions
associated with the Protestant reformations, not modern comparative theology;
opponents require their own voices and settings.

The printed \emph{Disputationes de controversiis christianae fidei} emerged
from this work in stages at Ingolstadt from 1586 through 1593.  Lectures and
books are not identical evidence.  Printing changed audience and arrangement;
successive impressions and later correction show revision rather than one
timeless summa dictated in finished form.  Bellarmine's 1608
\emph{Recognitio librorum omnium} is itself evidence that he revisited his
printed corpus and contested errors of production and interpretation.

\subsection{Authority and its limits}

The \emph{Disputationes} established Bellarmine's European reputation because
it joined scale to usability.  It made the divisions of western Christianity
look capable of ordered treatment and gave Catholic teachers an extensive
archive of replies.  Its historical argument could recover patristic and
conciliar evidence neglected in a narrower textbook.  It could also place
complex movements under categories selected by a confessional adversary.
Source abundance did not remove the need to assess transmission, attribution,
translation, or the fairness of representation.

One politically consequential distinction concerned papal power.  Bellarmine
denied that the pope possessed a direct universal dominion over temporal
affairs simply as temporal affairs.  He nevertheless defended an indirect
power to intervene in them when ordered to spiritual ends.  The distinction
limited one theory of direct control but retained a strong confessional claim
over rulers and communities.  It is not an early statement of modern
church--state separation, popular sovereignty, or civil religious liberty.

The same corpus contains a harder limit.  In the 1587 Ingolstadt printing of
\emph{De laicis} III.21, at printed page 495, Bellarmine defended temporal
punishment, including death, for incorrigible heretics, especially those who
had relapsed.  That
determination belongs to his theory of Christian society and cannot be moved
into an opponent's column.  It also cannot prove his personal vote in every
inquisitorial case.  Reading author, chapter, institution, and particular
dossier together prevents both moral evasion and unsupported assignment of
individual acts.

\evidenceaxis{The \emph{Disputationes} grew from Bellarmine's Roman College
lectures and appeared in a staged, revised print history.}{The autobiography,
early editions, the 1608 \emph{Recognitio}, and modern bibliographic
reconstruction control the sequence.}{No one printing is treated as every
lecture or final intention; a Catholic controversial work supplies its own
arguments, not a neutral account of every opponent.}
